\section{Архитектура системы и подход к реализации}
\label{sec:architecture}

Обзор литературы показал, что ни одна из существующих платформ не реализует
полный цикл от аналитики ученика к персонализированному конспекту с
математическими визуализациями. В настоящем разделе описывается, как этот
цикл устроен в разрабатываемой системе: какова её общая архитектура, как
устроены её модули и какие инженерные решения положены в их основу.
Ограничения текущей реализации и перспективные направления развития
рассматриваются отдельно в подразделе~\ref{sec:ch3:roadmap}.

\subsection{Общая архитектура системы}

Разрабатываемый модуль встраивается в платформу <<Прорыв>>~\cite{Proriv}
в качестве компонента, отвечающего за генерацию учебных материалов.
На входе модуль получает вектор компетенций ученика, на выходе ---
структурированный конспект с текстом и визуальными иллюстрациями,
адаптированными под его уровень.

Вектор компетенций формируется по истории взаимодействий ученика с платформой:
из решённых задач, допущенных ошибок и затраченного времени строится оценка
уровня освоения по каждой теме. Тематическое пространство опирается на
кодификатор ФИПИ~\cite{FIPI2025Demo}, однако не совпадает с ним в точности:
темы кодификатора служат основой, которая при необходимости детализируется
или перегруппируется с учётом характерных паттернов ошибок учеников.
Такая привязка принципиальна: она соотносит профиль ученика с реальной
структурой экзамена. Вектор компетенций выполняет роль управляющего сигнала
на всех последующих этапах генерации: определяет, какие темы включить в
конспект, какой уровень детализации выбрать и где добавить дополнительные
визуальные пояснения.

Генерация конспекта организована как двухэтапный процесс. На первом этапе
языковая модель формирует текстовое содержание по фиксированной структуре
обязательных блоков, адаптируя каждый блок под профиль ученика. На втором
этапе отдельный модуль получает текстовые описания нужных иллюстраций и
строит по ним математические диаграммы. Оба этапа описаны в следующих
подразделах.

\subsection{Генерация текстового содержания конспекта}

Центральный элемент первого этапа --- управляемая генерация текста
(controllable generation)~\cite{Liang2024Controllable}.
Конспект строится по фиксированной структуре, включающей несколько обязательных
блоков: определения и формулировки, разбор типичных ошибок, примеры задач
с пошаговым решением и визуальные пояснения. Фиксированная структура
гарантирует, что конспект покрывает все необходимые элементы независимо от
темы и уровня ученика. Языковая модель адаптирует содержание каждого блока
под вектор компетенций: для ученика с низким освоением темы включаются
промежуточные шаги рассуждения и больше примеров, для более подготовленного
ученика изложение компактнее.

Архитектура модуля намеренно построена без привязки к конкретной языковой
модели: интерфейс взаимодействия с моделью унифицирован, что позволяет
подменять её без изменения остальных компонентов системы.

Фактическая точность содержания обеспечивается за счёт генерации с
извлечением информации (retrieval-augmented generation,
RAG)~\cite{Lewis2020RAG}. Корпус для поиска формируется из материалов
открытого банка заданий ФИПИ~\cite{FIPI2025Bank}, что привязывает содержание
конспекта к требованиям экзамена. На этапе генерации модель обращается к
извлечённым фрагментам, что снижает риск фактических ошибок: некорректное
определение или ошибочная формула в учебном материале воспринимается учеником
как авторитетное утверждение и может закрепить ложное
знание~\cite{HallucinationSurvey2023}.

Модуль текстового этапа формирует все обязательные блоки конспекта в
соответствии с фиксированной структурой и определяет места для визуальных
элементов, передаваемых на следующий этап.

\subsection{Генерация визуализаций}

Второй этап отвечает за автоматическое создание математических иллюстраций.
Прямое использование моделей генерации изображений по тексту (text-to-image)
для этой задачи неприемлемо: такие модели систематически нарушают точные
геометрические отношения и некорректно воспроизводят подписи и математические
формулы~\cite{Bosheah2025,Wang2025MagicGeo}. В данной работе принят
альтернативный подход: языковая модель формирует не изображение, а
структурированную спецификацию визуального элемента, по которой
специализированный движок строит итоговое изображение.

Процесс состоит из трёх шагов. На первом шаге языковая модель, опираясь на
текстовое описание нужной иллюстрации, генерирует JSON"=спецификацию
визуального элемента. Формат спецификации разработан в рамках данной работы
и содержит семантически значимые поля: тип диаграммы, перечень объектов с их
параметрами, подписи и пространственные ограничения. На втором шаге
корректность выходного JSON обеспечивается библиотекой Instructor, которая
валидирует результат по Pydantic-схеме и при необходимости автоматически
повторяет запрос. На третьем шаге специализированный
движок разбирает спецификацию и строит изображение в формате масштабируемой
векторной графики (Scalable Vector Graphics, SVG)~\cite{SVG2Spec}.
SVG выбран потому, что он обеспечивает геометрическую корректность без потерь
при масштабировании, поддаётся программной проверке и может быть встроен
непосредственно в веб-интерфейс платформы.

Как и в модуле генерации текста, архитектура не зависит от конкретной языковой
модели: замена модели не требует изменений в схеме JSON или логике рендеринга.

Реализованный движок поддерживает три типа математических иллюстраций.

\textbf{Графики функций.} Движок строит графики элементарных и составных
функций, встречающихся в задачах ЕГЭ: степенных, показательных,
логарифмических, тригонометрических. Все значения функции вычисляются
программно по аналитическому выражению, что исключает погрешности ручного
построения. Поддерживаются подписи осей, отметки характерных точек (нули,
экстремумы, точки пересечения).

\textbf{Геометрические чертежи.} Для задач планиметрии и стереометрии
движок строит фигуры по заданным параметрам с корректными пространственными
соотношениями. Корректность обеспечивается на уровне спецификации:
JSON"=формат содержит явные ограничения (параллельность, перпендикулярность,
равенство отрезков), которые движок соблюдает при построении. Этот подход
согласуется с выводами Wang et al.~\cite{Wang2025MagicGeo}, показавших, что
корректные пространственные отношения требуют формальных ограничений,
поскольку свободная генерация систематически их нарушает. Также есть возможность создавать 
дополнительные построения, указывая только их смысл, например <<медиана из угла B>>, а координаты вычисляются программно. 

\textbf{Блок-схемы.} Для алгоритмических задач и объяснения многошаговых
методов решения движок строит блок"=схемы с соединительными стрелками и
условными ветвлениями. Есть возможность добавлять подписи к стрелкам. Также реализована возможность отрисовывать 
формулы и математические символы в формате LaTeX.

Этап генерации визуализаций реализован для всех трёх типов иллюстраций.
Расширение набора поддерживаемых геометрических примитивов и работа над
качеством подписей при сложных конфигурациях выполнены в виде
последовательных итераций улучшений и подробно описаны в
подразделе~\ref{sec:ch3:improvements}.

\subsection{Сценарии использования и контур развёртывания}

Разработанный модуль предназначен для использования в составе сервиса
<<Прорыв>> в качестве компонента генерации учебных материалов. Целевой
пользователь --- школьник, самостоятельно готовящийся к ЕГЭ по математике.

Связь между аналитическим профилем и генератором осуществляется через вектор
компетенций: каждое обращение к модулю сопровождается актуальным состоянием
профиля, что позволяет корректировать содержание конспекта при изменении
уровня ученика. Если ученик систематически допускает ошибки при исследовании
функции на монотонность, сгенерированный конспект уделяет этому пункту
больше внимания и включает дополнительный иллюстративный график.

Модуль реализован как самостоятельный компонент с унифицированным интерфейсом
обращения к языковой модели и фиксированным форматом входных и выходных
артефактов. Такой контур упрощает воспроизводимость экспериментов: прогон
системы не зависит от производственной инфраструктуры платформы и допускает
автономный запуск на отдельной рабочей машине. Сопряжение модуля с
действующим сервисом относится к эксплуатационным задачам и рассматривается
среди направлений развития (подраздел~\ref{sec:ch3:roadmap}).

\subsection{Оценка качества}

Оценка качества прототипа в рамках настоящей работы построена на двух
уровнях: техническая валидация процесса генерации и качественный анализ
полученных конспектов и иллюстраций. Принятый объём оценки соответствует
масштабу прототипа. На прототипном корпусе содержательные выводы о
работоспособности контура и о направлениях его улучшения получаются именно
из автоматических метрик и предметного разбора выходов, тогда как
полноценная проверка образовательного эффекта на реальных учениках
имеет смысл уже после первичной отладки и интеграции системы в действующий
сервис, и в прототипной фазе настоящей работы не проводилась. Обоснование
такого решения и возможные методические подходы к проверке на пользователях
обсуждаются среди направлений развития (подраздел~\ref{sec:ch3:roadmap}).

\textbf{Техническая валидация.} Технические свойства конвейера измерены
автоматическими метриками на корпусе из четырёх учебных конспектов
(вероятность, уравнения и неравенства, производная и её приложения,
прикладная математика) и восьми отдельных задач планиметрии и стереометрии.
По конспектной ветке проведена серия из трёх независимых прогонов с
фиксацией числа запросов, объёма входных и выходных токенов, средней
задержки генерации и числа запланированных визуалов. Результаты сведены в
табл.~\ref{tab:ch3:run-metrics} в подразделе~\ref{sec:ch3:integration}.
Дополнительно выполнены температурная абляция планировщика
(табл.~\ref{tab:ch3:temperature-ablation}) и сопоставление двух языковых
моделей разной размерности на одинаковом наборе входов
(табл.~\ref{tab:ch3:model-comparison}). Геометрическая ветка получила
симметричное журналирование на уровне отдельных шагов: на корпусе из восьми
задач проведена серия повторных прогонов с агрегированием времени, числа
токенов, доли планов, потребовавших исправления (fix\nobreakdash-rate), и
доли прогонов с визуализированным ответом. Результаты сведены в
табл.~\ref{tab:ch3:geometry-variance}. Эффект включения
пост\nobreakdash-валидатора визуализации ответа в режим замкнутого цикла
исправления зафиксирован отдельным экспериментом
(табл.~\ref{tab:ch3:answer-vis-before-after}). Корректность вычислений
агрегаторов и пост\nobreakdash-валидатора покрыта изолированными тестами
модулей \texttt{pipeline/\_smoke\_geometry\_variance.py} и
\texttt{pipeline/\_test\_answer\_validator.py}, не требующими обращения к
языковой модели.

\textbf{Качественный анализ конспектов и иллюстраций.} Сгенерированные
конспекты и геометрические чертежи проанализированы авторами работы на том
же корпусе. Этот анализ выявил характерные классы технических сбоев
(ошибки разбора JSON, нарушения ссылочной целостности в плане, коллизии
подписей, обрезка точек границами области, ошибки порядка обхода вершин в
трёхмерных гранях, пропуск визуального представления ответа) и характерные
педагогические дефекты, систематически описанные в
подразделе~\ref{sec:ch3:limits}. Разбор стал содержательной базой для
тринадцати итераций улучшений (подраздел~\ref{sec:ch3:improvements}):
каждая итерация привязана к конкретному классу наблюдавшихся ошибок и
переносит часть рассуждения с языковой модели на детерминированные
алгоритмы. На прототипном масштабе именно такой предметный разбор
обеспечивает содержательную валидацию инженерных решений и позволяет
поставить осмысленные вопросы для последующего экспертного и полевого
исследований.

Сочетание двух описанных уровней --- технического и качественного ---
достаточно для того, чтобы зафиксировать работоспособность контура и
обоснованность принятых инженерных решений на прототипном корпусе. Более
развёрнутые формы валидации (формальное сравнение с базовой линией
text\nobreakdash-to\nobreakdash-image, экспертная разметка иллюстраций
учителями математики, контролируемое тестирование на учениках) требуют
отдельных протоколов, экспертного аппарата и доступа к производственному
окружению. Их обоснование, возможные методические подходы и место в плане
работ обсуждаются среди направлений развития системы
(подраздел~\ref{sec:ch3:roadmap}).

\paragraph{Вклад авторов в данную главу.} Описание общей архитектуры системы и модуля генерации текстового содержания конспекта подготовлено А.~С.~Низовым. Описание модуля генерации визуализаций, сценариев использования и подхода к оценке качества подготовлено Д.~А.~Ковырковым. Двухэтапная декомпозиция процесса и интерфейс между модулями согласованы авторами совместно.